\section{The locus of Hodge classes}\label{app:cdk}

Deformation arguments about Hodge classes rest on one theorem, that of
Cattani, Deligne and Kaplan \cite{CDK95}: the locus in a family where a given
class stays a Hodge class is an algebraic subvariety, not merely an analytic
one. This appendix states the theorem, proves its elementary part, isolates
the two analytic inputs that the rest requires, computes the example that
matters here, and explains why the theorem is not a route to the Hodge
conjecture. The last point is the one that bears on \Cref{sec:construction},
and it is why the material is set out here rather than cited.

\subsection{Variations of Hodge structure}

\begin{definition}[Variation of Hodge structure]\label{def:vhs}
Let $S$ be a smooth connected complex quasi-projective variety. A
\emph{polarised variation of $\QQ$-Hodge structure} of weight $w$ on $S$
consists of
\begin{enumerate}[label=\textup{(\alph*)},nosep]
\item a local system $\mathbb{V}$ of finite dimensional $\QQ$-vector spaces on
  $S$, with associated holomorphic bundle
  $\cV=\mathbb{V}\otimes_{\QQ}\cO_{S}$ carrying the flat connection $\nabla$
  whose flat sections are $\mathbb{V}$;
\item a decreasing filtration $\cV=F^{0}\supseteq F^{1}\supseteq\cdots$ by
  holomorphic subbundles, inducing on each fibre $\mathbb{V}_{s}$ a
  $\QQ$-Hodge structure of weight $w$;
\item a flat nondegenerate $(-1)^{w}$-symmetric pairing
  $Q\colon\mathbb{V}\otimes\mathbb{V}\to\QQ_{S}$ polarising each fibre;
\end{enumerate}
subject to Griffiths transversality
\begin{equation}\label{eq:transversality}
  \nabla F^{p}\subseteq F^{p-1}\otimes\Omega^{1}_{S}.
\end{equation}
\end{definition}

The example to keep in mind is the geometric one: for $f\colon\cX\to S$ smooth
projective, $\mathbb{V}=R^{w}f_{*}\QQ$ with its Gauss--Manin connection is such
a variation, and \eqref{eq:transversality} is Griffiths' theorem
\cite[Ch.~10]{Voi02}. The other example we use is the Weil family of
\Cref{prop:weildomain} below, which is defined in the first instance by a
period domain rather than by a family of varieties, and which carries a
variation all the same.

\begin{notation}\label{not:totalspace}
Write $\cV^{\mathrm{an}}$ for the total space of $\cV$, a complex manifold
fibred over $S$ with fibre $\mathbb{V}_{s}\otimes\CC$, and
$\mathbb{V}^{\mathrm{tot}}\subset\cV^{\mathrm{an}}$ for the subset of pairs
$(s,\alpha)$ with $\alpha\in\mathbb{V}_{s}$, a covering space of $S$ with
countable fibres.
\end{notation}

\subsection{Hodge loci}

\begin{definition}\label{def:hodgelocus}
Let $\mathbb{V}$ be a polarised variation of weight $2p$ on $S$. Its
\emph{locus of Hodge classes} is
\[
  \Hdg^{p}(\mathbb{V})
  =\bigl\{(s,\alpha)\in\mathbb{V}^{\mathrm{tot}}\;:\;
     \alpha\in F^{p}\cV_{s}\bigr\}.
\]
For $N\in\QQ$ write $\Hdg^{p}_{N}$ for the part with $Q(\alpha,\alpha)=N$.
\end{definition}

By conjugation symmetry, a rational $\alpha$ of weight $2p$ lies in $F^{p}$
exactly when it is of type $(p,p)$. The following two lemmas are elementary,
and they give everything about $\Hdg^{p}$ except algebraicity.

\begin{lemma}[Analyticity]\label{lem:hodgeanalytic}
$\Hdg^{p}(\mathbb{V})$ is a closed analytic subset of
$\mathbb{V}^{\mathrm{tot}}$, and the projection
$\pi\colon\Hdg^{p}(\mathbb{V})\to S$ is locally biholomorphic onto its image on
each connected component. In particular each component is a closed analytic
subvariety of an open subset of $S$, and there are countably many components.
\end{lemma}

\begin{proof}
Work on a simply connected open $U\subseteq S$ over which $\mathbb{V}$ is
trivialised, $\mathbb{V}|_{U}\cong V_{0}\times U$ for a fixed $\QQ$-vector
space $V_{0}$. Then $\mathbb{V}^{\mathrm{tot}}|_{U}$ is the disjoint union over
$\alpha\in V_{0}$ of copies of $U$. Fix $\alpha$. The quotient
$\cV/F^{p}$ is a holomorphic vector bundle on $U$, and the image
$\bar\alpha(s)$ of the flat section $\alpha$ in $(\cV/F^{p})_{s}$ is a
holomorphic section of it. The set of $s\in U$ with $\alpha\in F^{p}\cV_{s}$
is the zero locus of $\bar\alpha$, a closed analytic subset of $U$. Taking
the union over the countable set $V_{0}$ gives the first statement, and shows
that on the sheet indexed by $\alpha$ the projection $\pi$ is the inclusion of
that zero locus into $U$, hence injective. There are countably many
components because $S$ is covered by countably many such $U$, the set $V_{0}$
is countable, and each zero locus has countably many components.
\end{proof}

\begin{lemma}[Finiteness at bounded norm]\label{lem:boundednorm}
Fix $s\in S$ and $N\in\QQ$. The set of $\alpha\in\mathbb{V}_{s}$ with
$\alpha$ of type $(p,p)$, $Q(\alpha,\alpha)=N$, and $\alpha$ in a fixed
lattice $\mathbb{V}_{s,\ZZ}$, is finite. The number of such $\alpha$ is locally
constant on each component of $\Hdg^{p}_{N}$.
\end{lemma}

\begin{proof}
Decompose $\alpha$ into its primitive parts for the Lefschetz decomposition
attached to the polarisation. On the primitive part of type $(p,p)$ in weight
$2p$ the Hodge--Riemann bilinear relations read
$(-1)^{p}\,Q(\alpha,\bbar{\alpha})>0$ for $\alpha\ne0$, so the form
$(-1)^{p}Q$ is positive definite on the real span of the primitive real
$(p,p)$-classes; the same holds on each Lefschetz summand with an alternating
sign, and the resulting form $\tilde Q$ obtained by flipping signs summand by
summand is positive definite on the whole space of real $(p,p)$-classes. The
conditions $Q(\alpha,\alpha)=N$ and $\alpha$ of type $(p,p)$ therefore confine
$\alpha$ to a compact set for $\tilde Q$, and a compact set meets a lattice in
finitely many points. The count is locally constant because the lattice and $Q$
are flat and the Hodge filtration varies continuously, so no Hodge class can
appear or disappear without its $Q$-norm changing.
\end{proof}

\begin{theorem}[Cattani, Deligne and Kaplan \cite{CDK95}]\label{thm:cdk}%
\footnote{The paper is short and the proof is not: the two inputs isolated
below, the nilpotent orbit theorem and the $\SL_{2}$-orbit theorem, carry the
weight.}
Let $S$ be a smooth complex quasi-projective variety and $\mathbb{V}$ a
polarised variation of $\ZZ$-Hodge structure of weight $2p$ on $S$. Then each
connected component of $\Hdg^{p}(\mathbb{V})$ is a closed algebraic subvariety
of $\mathbb{V}^{\mathrm{tot}}$, and its image in $S$ is a closed algebraic
subvariety of $S$.
\end{theorem}

We do not reproduce the proof, which draws on the whole asymptotic theory of
degenerating Hodge structures developed in the two preceding decades. For
families of abelian varieties the algebraicity of these loci also follows, by
a different route, from Deligne's theorem that Hodge classes on abelian
varieties are absolutely Hodge; \Cref{thm:cdk} is what carries it beyond
abelian varieties. We record where the difficulty lies and which two
theorems remove it, since the shape of the argument is what one needs in
order to judge whether a deformation argument in this area is legitimate.

By \Cref{lem:hodgeanalytic} a component $Z$ of $\Hdg^{p}$ is already a closed
analytic subvariety of an open subset of $S$, and by \Cref{lem:boundednorm} it
carries a locally constant $Q$-norm and is one sheet of a finite cover of its
image. Choose a smooth projective compactification $\bar S\supseteq S$ with
$\bar S\setminus S$ a normal crossing divisor. By Chow's theorem it is enough
to show that $Z$ extends to an analytic subvariety of $\bar S$, that is, that
$Z$ does not oscillate wildly as one approaches the boundary. This is an
asymptotic question about the Hodge filtration, and it is answered by two
theorems.

\begin{enumerate}[label=\textup{(\Alph*)},leftmargin=2.2em]
\item \emph{The nilpotent orbit theorem} of Schmid \cite{Sch73}: near a
  boundary point the period map is approximated, to exponentially small error
  in the natural coordinates, by the orbit of a nilpotent Lie algebra element
  built from the local monodromy logarithms.
\item \emph{The $\SL_{2}$-orbit theorem} in several variables, of Cattani,
  Kaplan and Schmid \cite{CKS86}: that nilpotent orbit is in turn approximated
  by an orbit of a fixed real reductive group, with uniform estimates in the
  several-variable setting.
\end{enumerate}

Together these give a two-sided estimate for the Hodge norm of a flat section
near the boundary, of the shape $\|\alpha\|_{s}\sim\prod_{j}(\log|t_{j}|^{-1})^{m_{j}}$
with integers $m_{j}$ determined by the weight filtration. A flat section that
is Hodge on $Z$ has its norm for the positive definite form $\tilde Q$ of
the proof of \Cref{lem:boundednorm} bounded by that lemma, and the
estimate converts that bound into a statement that the possible boundary
behaviours form a finite list. Algebraicity then follows from the analytic
extension across the boundary together with Chow's theorem in $\bar S$.

\begin{remark}\label{rem:cdkinputs}
Both (A) and (B) are theorems about degenerating Hodge structures, and
neither mentions algebraic cycles. The algebraicity in \Cref{thm:cdk} is
algebraicity of a locus of periods, obtained by analysis, and it carries no
information about algebraic cycles; \Cref{ssec:cdknot} develops the point.
\end{remark}

\subsection{The Weil family}

The example relevant to this paper is the family of complex structures on the
model of \Cref{setup:model}.

\begin{proposition}[The Weil period domain]\label{prop:weildomain}
Fix $K$, $n$, the $K$-space $V=K^{2n}$, and the hermitian form $H$ of
signature $(n,n)$ from \Cref{lem:hermform}. The set of Hodge structures of
weight one on $V$ that are polarised by $Q$, commute with the $K$-action, and
are of $(K,-1,n)$-Weil type is in bijection with
\[
  \cD_{n,n}
  =\bigl\{P\subset V_{+}\;:\;\dim_{\CC}P=n,\ H|_{P}\ \text{negative definite}\bigr\},
\]
an open subset of the Grassmannian $\mathrm{Gr}(n,2n)$ and a bounded symmetric
domain isomorphic to $\mathrm{U}(n,n)/\bigl(\mathrm{U}(n)\times\mathrm{U}(n)\bigr)$.
In particular
\[
  \dim_{\CC}\cD_{n,n}=n^{2},
\]
compared with $\dim\cA_{2n}=n(2n+1)$ for the moduli of principally polarised
abelian $2n$-folds.
\end{proposition}

\begin{proof}
A Hodge structure of weight one on $V$ is a subspace $V^{1,0}\subset V_{\CC}$
with $V_{\CC}=V^{1,0}\oplus\overline{V^{1,0}}$. Commuting with the $K$-action
means that $V^{1,0}$ is stable under the operator $S$ of \Cref{setup:model},
so that it splits as $V^{1,0}=P\oplus P'$ with $P=V^{1,0}\cap V_{+}$ and
$P'=V^{1,0}\cap V_{-}$.
Conjugation interchanges $V_{+}$ and $V_{-}$, so
$\overline{V^{1,0}}=\bbar{P}\oplus\bbar{P'}$ with $\bbar{P}\subset V_{-}$, and
the direct sum condition becomes $V_{-}=P'\oplus\bbar{P}$, which determines
$P'$ up to the choice of a complement; the Weil condition
$\dim P=\dim P'=n$ of \Cref{lem:balanced} together with the polarisation
condition pins it down as $P'=\bbar{P}^{\perp_{H}}$. The structure is
therefore determined by $P$.

For the positivity, extend $H$ to the $\CC$-bilinear pairing between $V_{+}$
and $V_{-}$. By the computation \eqref{eq:Qplusminus} in the proof of
\Cref{thm:model}, the second Riemann relation holds exactly when $i\,Q$ is
positive on $V^{1,0}$, and on the $V_{+}$ part this is the condition that $H$
be negative definite on $P$, the sign being fixed by the convention of
\Cref{lem:hermform}. The set of $n$-dimensional negative definite subspaces
of a hermitian space of signature $(n,n)$ is the standard realisation of the
symmetric domain of $\mathrm{U}(n,n)$: the unitary group acts transitively on
it and the stabiliser of a point is $\mathrm{U}(n)\times\mathrm{U}(n)$, the
unitary groups of $P$ and of $P^{\perp}$. It is open in $\mathrm{Gr}(n,2n)$
because negative definiteness is an open condition, and
$\dim\mathrm{Gr}(n,2n)=n(2n-n)=n^{2}$.

For the comparison, $\dim\cA_{g}=g(g+1)/2$, which is $n(2n+1)$ for $g=2n$.
\end{proof}

\begin{corollary}[The Weil locus is algebraic]\label{cor:weilalgebraic}
Inside $\cA_{2n}$ the image of $\cD_{n,n}$, that is the locus of principally
polarised abelian $2n$-folds admitting an action of an order of $K$ of
$(K,-1,n)$-Weil type, is a closed algebraic subvariety of dimension $n^{2}$.
\end{corollary}

\begin{proof}
The condition of admitting such an action is the condition that a certain
class in $\End(H^{1})$, namely the graph of $\iota(\sqrt{-d})$, be a Hodge
class, together with the balanced signature, which is locally constant. The
locus is therefore a component of a locus of Hodge classes for the variation
$\End(R^{1}f_{*}\QQ)$ over $\cA_{2n}$ with its level structure, and
\Cref{thm:cdk} applies. The dimension is \Cref{prop:weildomain}.
\end{proof}

\subsection{Limits of the theorem}\label{ssec:cdknot}

\begin{remark}[Algebraic locus, not algebraic class]\label{rem:notalgebraic}
\Cref{thm:cdk} says that the set of parameters where $\alpha$ is Hodge is cut
out by polynomial equations. It says nothing about whether $\alpha$ is the
class of a cycle at any one of those parameters. The two statements are
independent in both directions: the Hodge locus of a class that happens to be
algebraic everywhere is still known to be algebraic only through
\Cref{thm:cdk}, and conversely the algebraicity of the locus supplies no
cycle. The distinction is easily blurred, and \Cref{cor:weilalgebraic} is
exactly the kind of statement that invites it: the Weil locus is an algebraic
subvariety, even a Shimura subvariety, and none of that produces the class
$\omega_{1}$ of \eqref{eq:omegagens}. \Cref{fig:moduli} compares the
dimension of the Weil locus with that of $\cA_{2n}$.
\end{remark}

\begin{figure}[htbp]
\centering
  \includegraphics[width=0.92\textwidth]{fig_moduli}
\caption{The Weil locus inside the moduli of principally polarised abelian
$2n$-folds, for $n=1,\dots,6$. Each prism has the height $n(2n+1)$ of the
moduli space $\cA_{2n}$; its lower part, of height $n^{2}$, is the
dimension of the period domain $\cD_{n,n}$ of \Cref{prop:weildomain}, and
its upper part is the codimension $n(n+1)$. The codimension grows
quadratically, so the Weil locus is thin, and by \Cref{cor:weilalgebraic} it
is an algebraic subvariety. Both facts matter for propagation: the locus
along which one would have to deform is small and is a variety, but neither
fact produces a cycle on it.}
\label{fig:moduli}
\end{figure}

\begin{remark}[The variational Hodge conjecture]\label{rem:variational}
The property that would turn \Cref{thm:cdk} into a tool is the variational
Hodge conjecture of Grothendieck: if $\cX\to S$ is smooth projective over a
connected base, $\alpha$ is a flat section of $R^{2p}f_{*}\QQ$ that is Hodge at
every point, and $\alpha_{s_{0}}$ is algebraic at one point $s_{0}$, then
$\alpha_{s}$ is algebraic at every point. Its infinitesimal form is a
theorem under a semiregularity hypothesis: if the cycle at
$s_{0}$ is semiregular in the sense of Bloch \cite{Blo72}, then it deforms to
first order, and hence formally, along the locus where its class stays Hodge.
That theorem is the subject of \Cref{sec:semireg}, and
\Cref{cor:countnotbind} shows that its numerical hypothesis is not what
obstructs the present problem.
\end{remark}

\begin{remark}[The construction of
\texorpdfstring{\Cref{sec:construction}}{Section 14}]\label{rem:cdkconstruction}
The construction of \Cref{sec:construction} produces, for every $n$, every $d$
and every nondegenerate $\beta$, an abelian variety of Weil type together with
secant sheaves whose Chern classes are computed in closed form, and its output
is algebraic on the split locus. The missing ingredient, described in
\Cref{rem:deformation}, is a statement along the algebraic locus of
\Cref{cor:weilalgebraic}: the Hodge type of the deformed object has to be
controlled as the parameter moves off the split locus. \Cref{thm:cdk}
guarantees that the locus along which one deforms is an algebraic variety,
which is what makes the question well posed, but it does not answer the
question. \Cref{q:higher} states it precisely.
\end{remark}
